% Kapitel 10 -- Von Kuonr\^{a}ts Streit vor dem Rat zu Znaim

\section{Von Kuonr\^{a}ts Streit vor dem Rat zu Znaim}

\mylettrine{A}{ls} Kuonr\^{a}t von Steinare in jenen Tagen nach \gls{Znaim} ritt, war die Stadt in großer Unruhe. Man hatte den alten \gls{Gruober} gefangen, so heisst es, und sein Zutun ward in scharfen Worten als Schandtat kundgetan. \gls{Kuonrat von Steinare}, wenig bedacht auf List und Rücksicht, verlangte vor dem Rat die Freiheit des Gefangenen.

Er trat vor die Rathausstufen, und man lachte ihn an, wie man eben lacht, wenn ein Vagabund die Stadttore vernimmt und denkt, ihm gelte ebensoviel Recht wie einem Herrn. Doch \gls{Kuonrat von Steinare} liess sich nicht schelten. In einer langen und konfusen Rede, wie mir von verschiedenen Gewährsmännern berichtet ward, sprach er mit den Worten eines Mannes, der seine Sache mit dem Herzen, nicht mit dem Gesetz und nicht mit dem Verstand, vorträgt. Ich gebe hier das Wesentliche jener Ansprache wieder, so gut es die Zeugen erinnern konnten, wiewohl keiner von ihnen des Fremden Gedankengang bis zum Ende zu verfolgen vermochte. Er sagte:

\enquote{Es ist gegen die \gls{Freiheit}, einen guten Geschäftsmann gefangen zu nehmen. Ihr wisst das vielleicht nicht, denn ihr sitzt hier hinter euren Mauern, aber dieser Mann ist kein Bösewicht. Er ist ein Handwerker, der eine Arbeit verrichtet, die sonst niemand verrichtet, und die doch getan werden muss. Denn es gibt Frauen in diesen Landen, viele Frauen, die hungern, und ich frage euch: was ist schlimmer, eine Frau, die hungert, oder eine Frau, der man Dienst anbietet? Ich warte. Ihr antwortet nicht, weil ihr die Antwort kennt. Schlimmer ist die hungernde, das wird jeder verständige Mann zugeben. Also ist dieser Mann, indem er Frauen Dienst angeboten hat, ein Wohltäter. Versteht ihr das? Ein Wohltäter. Und die Frauen, die mit ihm gegangen sind, die sind freiwillig gegangen, niemand hat sie gezwungen, das ist ihr gutes Recht, denn die \gls{Freiheit} ist das höchste Gut, das Gott dem Menschen gegeben hat, und wer diese \gls{Freiheit} beschneidet, der ist der Feind, nicht der Mann, der den Frauen Arbeit gibt. Ich sage euch das klar und deutlich. Nun, ich habe mit dieser Sache selbst nur wenig zu tun, das sei klar gesagt, denn ich bin nur hergeritten, weil ich diesen Mann kenne und weil mir unrecht schien, was ich hörte. Mein Ross \gls{Veltloufer} war müde von der Fahrt, aber das ist nicht wichtig, ich sage es nur, damit ihr wisst, dass ich es nicht leichtfertig nehme. Also, wo war ich. Ja: die Frauen sind freiwillig mitgegangen. Und dieser Mann ist ein Geschäftsmann. Ihr sperrt ihn ein, und damit schadet ihr dem mährischen Volk, denn wer gibt diesen Frauen jetzt Arbeit, wenn der Mann, der sie vermittelte, hinter Schloss und Riegel sitzt? Sagt mir das. Ihr werdet es nicht sagen können, denn es gibt keine Antwort. Ich habe auch gehört, dass man in dieser Stadt Händler duldet, die schlechtes Getreide verkaufen, und denen tut ihr nichts. Ob das stimmt, weiss ich nicht genau, aber es wäre nicht verwunderlich. Ich will damit sagen: ihr seid nicht gerecht in eurer Strenge, und wenn ihr meint, ich irre mich, so sagt mir das, und ich werde zuhören, denn ich bin kein unverständiger Mann, das ist bekannt. Aber dieser gute Geschäftsmann muss frei sein, denn er hat Frauen Arbeit gegeben, und das ist kein Verbrechen, sondern ein Dienst an Gott und an den Menschen, und eure Entscheidung schadet dem mährischen Volk und der \gls{Freiheit} insgesamt, und das sage ich mit aller Bestimmtheit.}

Hier, so berichten die Zeugen, hielt Kuonr\^{a}t inne und schaute den Rat erwartungsvoll an, als habe er eben einen unwidersprechlichen Beweis vorgetragen. Die Ratsherren sahen einander an. Mancher wusste nicht mehr, was der Fremde eigentlich begehre; ob er den Gefangenen freikaufen wolle, ob er ein Recht auf Auskunft erhebe, oder ob er schlicht die gesamte Ordnung der Stadt in Frage stelle. Einige berichteten später, er habe dieselbe Sache dreimal in wechselnden Worten vorgetragen, und dazwischen ungebeten von der Müdigkeit seines Pferdes geredet und von der Ungerechtigkeit der Welt im Allgemeinen. Andere meinten, er habe sich selbst nicht sicher gewusst, wieviel er mit der Sache zu tun habe. Der Herr allein weiss, was \gls{Kuonrat von Steinare} in jener Stunde wirklich glaubte.

Die Herren des Rates blieben ungerührt und erwiderten mit ruhigem Zorn, wie es ihres Amtes haftet. Einer unter ihnen, ein älterer Mann mit grauem Bart, sprach mit schneidender Kürze: \enquote{Dies ist kein ehrlicher Handel, sondern Raub und Verkauf von Menschen. Wir kennen die Nachrichten von entführten Mägden. Man bringt sie fort, um sie weiterzuverkaufen. Dergleichen darf nicht geduldet werden, denn es ist nicht nur ein Vergehen gegen unser weltlich Recht, es ist auch ein Frevel gegen Gottes Gebot.}

Auf diese Worte hin geriet \gls{Kuonrat von Steinare} vollends außer sich. Er schrie, so heisst es, dass die Tauben vom Dach flogen, und rief: \enquote{Ihr seid Heuchler, alle miteinander! Ihr sitzt hier wie Fürsten, aber ihr seid Diener von Bischöfen und Herren, nichts weiter! Wer hat euch das Recht gegeben, einen freien Mann vor euch herzuzitieren? Dieser gute Mann, den ihr da eingekerkert habt, der hat den Frauen gedient, und das ist mehr, als ihr in eurem ganzen Leben getan habt! Und ich sage noch mehr. Ich sage sehr viel mehr. Ihr kennt die \gls{Freiheit} nicht, ihr habt sie nie gekannt, denn wer hinter Mauern hockt und auf andere zeigt, der ist selbst nicht frei!} Es ward noch lauter, und man verstand, so berichten die Zeugen, im einzelnen nicht mehr alles, was er rief, denn Zorn und Verwirrung liefen in seiner Rede ineinander.

Da der Rat seine Ruhe nicht wiederfand und die Stadtväter fürchteten, ein Aufruhr möge folgen, liess man \gls{Kuonrat von Steinare} festnehmen. Man band ihm die Hände, man führte ihn ab und sperrte ihn in ein städtisches Haus, bis das verderbte Getümmel sich lege. Sein \gls{Knecht} hingegen blieb frei; er war es, der schweigend dem Herrn gefolgt war und der, wie berichtet wird, sich draußen vor den Toren herumdrückte, bis die Gefahr vorbei war.

Was wahr ist an manchem Wort der einen und der andern Seite, bleibt der Kunde für späterer Tage offen. So viel ist aber sicher, wie es die Stadtchroniken verzeichnen: \gls{Kuonrat von Steinare} ward niedergeführt, und die Stadt verwahrte ihn, bis man das Recht und die Sitte abgewogen hatte.

Währenddessen wurde der Spielmann und Räuber \gls{Gruober}, ganz wie es dem Gesetz und der Sitte entspricht, vor die Stadtmauern geführt und dort mit einem Strick um den Hals aufgehängt, damit alle sehen konnten, was mit solchen Leuten geschieht.
